\section{Conclusion}
\label{sec:conclusion}

The Census Bureau's campus-counting rule produces a systematic, directionally partisan redistricting effect: college towns are inflated by 10--35\% relative to what they would show under home-address counting, and that inflation benefits Democratic-leaning college-district configurations at the expense of the Republican-leaning suburban districts from which students originate.

We have documented this effect across the 50 largest college-town counties, with extended case studies for Ithaca, Ann Arbor, and Champaign-Urbana. Our estimate of approximately 180 tracts nationally that would change district assignment under home-address reallocation establishes the redistricting materiality of the issue. The effect is concentrated in competitive states---Michigan, Illinois, Pennsylvania, North Carolina, Wisconsin, and Iowa---where the partisan direction of campus counting is most consequential.

Legally, campus counting is permissible under \emph{Reynolds v. Sims} and \emph{Evenwel v. Abbott}: the Census Bureau's enumeration is entitled to deference, and states that use it for redistricting are complying with the one-person-one-vote standard. States have authority, however, to adopt home-address corrections for state redistricting purposes, as New York's prison-counting reform demonstrates. No state has yet done so for college students.

BISECT's \texttt{--population-source} flag enables systematic sensitivity analysis: a redistricting principal can directly compare bisect plans under campus-counted versus home-address-adjusted populations, observing which district boundaries shift and by how much. This transparency---seeing the redistricting consequence of a counting choice---is the principal contribution of algorithmic redistricting to the college-counting debate.

The college-student counting problem will not disappear. As college enrollment concentrations grow and as partisan sorting continues to align college towns with one political party, the redistricting consequences of campus counting will become more salient. This paper provides the analytical framework for that debate.

\paragraph{Acknowledgments.}
The author thanks the National Center for Education Statistics for IPEDS enrollment data and the Census Bureau for Group Quarters tabulations made publicly available under the 2020 Redistricting Data Program.
